\section{Complete cyclic overlap queries on compressed relators}
\label{sec:compressed-lcs}

\subsection{The local completeness gap being removed}

The preceding whole-donor fallback leaves a concrete gap: it can miss an
overlap longer than half the donor if that overlap uses only part of the
donor or starts at another cyclic position. The new fallback computes a
longest common substring of doubled relators, capped at one donor turn.
It therefore searches all cyclic rotations and both donor signs, without
expanding either word. It uses the existing ordinary relator certificate;
no new certificate format or trusted verifier operation is introduced.

The cut-pair and periodic-extension strategy follows Matsubara et
al.~\cite{matsubara-compressed-lcs}. Their implementation uses a particular
overlap/FM data structure and proves an $O(s^4\log s)$ arithmetic-operation
bound. Our implementation combines the existing deterministic SLP equality
engine with a dyadic construction of overlap progressions using the
Lifshits occurrence table. These substituted primitives have different
costs. We prove a polynomial local bound below; we do not assign the
published degree to this implementation. The later refinement in
Section~\ref{sec:lcs-bounds} adds stronger inexpensive bounds and
incremental witness processing; the measurements below describe the
initial complete-query checkpoint.

\subsection{Dyadic construction of suffix-prefix overlaps}

Write $X[i:j]$ for a zero-based half-open substring and define
\[
 \operatorname{OL}(X,Y)=
 \{k:1\le k\le\min(|X|,|Y|),\quad X[|X|-k:|X|]=Y[0:k]\}.
\]
We represent this set as disjoint arithmetic progressions of lengths.
For each power of two $l$, consider the band
$[l,h]$, where $h=\min(2l-1,|X|,|Y|)$. An overlap of length in this band
must begin with the seed $P=Y[0:l]$, occurring in
$X[|X|-h:|X|]$. The window has length less than $2l$; all seed occurrences
share a cut and are represented by an arithmetic progression. Their starts
$p$ correspond to candidate overlap lengths $|X|-p$.

A singleton candidate is checked by exact LCP. For a progression
$p_i=p_0+id$ with at least two members, $0<d<l$. Adjacent overlapping
occurrences imply that the seed has period $d$, and their union is one
$d$-periodic run. Compute its maximal right endpoint in $X$ and the maximal
$d$-periodic prefix length of $Y$:
\[
 E_X=p_0+d+\operatorname{LCP}(X[p_0:|X|],X[p_0+d:|X|]),
 \qquad
 E_Y=d+\operatorname{LCP}(Y,Y[d:|Y|]).
\]
Every candidate starts in the same periodic phase as $Y$.

\begin{lemma}[Filtering a seed progression]
If $E_X=|X|$, exactly the candidates with $p_i\ge |X|-E_Y$ are
suffix-prefix overlaps. If $E_X<|X|$, the only possible overlap candidate
is $p_i=E_X-E_Y$; when this value belongs to the progression, an exact
LCP query determines whether its entire suffix matches.
\end{lemma}
\begin{proof}
Compare the periodic run beginning at $p_i$ with the prefix of $Y$.
They agree until the smaller of the two periodic extents
$E_X-p_i$ and $E_Y$. If these extents differ and the shorter string has
not ended, the first departure from the periodic word causes a mismatch:
the longer side still has the periodic letter. Thus a match extending
through a nonperiodic tail of $X$ requires equal departure positions.
That equality determines at most one progression member. If $X$ ends
while still periodic, the whole suffix matches precisely when its length
is no greater than the periodic extent of $Y$. A possible end of $Y$
causes no exception, since every candidate length in the band is at most
$|Y|$.
\end{proof}

Clipping and reversing the surviving progression of starts gives a
progression of overlap lengths. Distinct bands are disjoint and cover
all positive lengths, so the resulting list is exact and has
$O(\log\min(|X|,|Y|))$ progressions. Uniform equal-letter words have the
direct answer: every positive length up to the shorter word length.
Different uniform letters have no overlaps. Neither case enumerates that set.

\subsection{From overlaps to longest common substrings}

Consider grammar rules $U=AB$ and $V=CD$. For each
$k\in\operatorname{OL}(A,D)$, let
\[
 L(k)=\operatorname{LCSuf}(A[0:|A|-k],C),\qquad
 R(k)=\operatorname{LCP}(B,D[k:|D|]).
\]
Then $k+L(k)+R(k)$ is a common substring length, with start positions
$|A|-k-L(k)$ in $U$ and $|C|-L(k)$ in $V$. Inversion reverses and
negates both words, so the existing LCP primitive on their inverses
computes LCSuf exactly. The alignment $k=0$ is also valid and uses
$\operatorname{LCSuf}(A,C)+\operatorname{LCP}(B,D)$.

\begin{lemma}[Coverage by pairs of grammar cuts]
Every common substring of length at least two is covered by one of these
alignments for some pair of binary rules, allowing $U$ and $V$ to be
interchanged. Single-letter matches can be handled separately.
\end{lemma}
\begin{proof}
In each occurrence's derivation tree, choose the lowest node containing
the whole occurrence. Since the occurrence has length at least two, it
crosses that node's binary cut. Place the two equal substrings on the same
coordinate line. If the first cut is to the right of the second by $k>0$,
the intervening word is both a suffix of $A$ and a prefix of $D$, and is
therefore in $\operatorname{OL}(A,D)$. Extending in both directions covers
the original common substring. Reverse the roles if the cuts are ordered
the other way; coincident cuts give $k=0$.
\end{proof}

Although the derivation trees may be exponential, each rule has a fixed
spelling. We examine rule pairs, not tree-node occurrences. A top-down
pass on the grammar DAG records the earliest position of each rule in
the root word. Adding that offset to a local witness gives a valid root
witness. We examine longer rules first and stop examining rules whose
length cannot improve the current best result.

\subsection{Only six members of a progression need extension}

Let $k=a+id$, $0\le i<t$, be an overlap progression with $t\ge2$.
The two largest overlaps imply that the common suffix of $A$ and prefix
of $D$ of length $a+(t-1)d$ have period $d$. Let $e_2$ be the maximal
$d$-periodic suffix length of $A$ and $e_3$ the maximal $d$-periodic
prefix length of $D$. Let $e_1$ be the extent to which $B$ continues the
periodic suffix of $A$, and $e_4$ the extent to which $C$ continues the
periodic prefix of $D$ backwards. All four values are obtained with
compressed suffix/prefix queries. For example,
\[
 e_2=d+\operatorname{LCSuf}(A[0:|A|-d],A[d:|A|]),
 \qquad e_3=d+\operatorname{LCP}(D,D[d:|D|]).
\]
For $e_1$, compare $B$ with enough binary-powered copies of the last
$d$ letters of $A$. For $e_4$, compare $C$ backwards with enough copies
of the first $d$ letters of $D$. Only a logarithmic number of new power
rules is required.

Except when periodic departure positions coincide, the extension length is
\[
 F(k)=k+\min(e_2-k,e_4)+\min(e_1,e_3-k).
\]
Its two breakpoints are $e_2-e_4$ and $e_3-e_1$. Between breakpoints it
is affine: first nondecreasing, then constant, then nonincreasing. At a
breakpoint, equal departures may match for longer; the true extension
must therefore be checked there. It follows that the maximum over the
progression is attained among its endpoints and the progression members
immediately below and above the two breakpoints, at most six values.
We evaluate those candidates by the exact extension formula, not by
assuming that the periodic approximation remains valid at a breakpoint.
Singleton progressions require one exact extension. This argument permits
arbitrary starting residues and does not require the first overlap to
be less than the step.

\subsection{Cyclic search, early bounds and certificate replay}

For relators $R,D$, every cyclic substring of length at most
$\min(|R|,|D|)$ appears in $RR$ and $DD$. Conversely, any such common
substring of the doubled words determines valid cyclic starting positions
by reducing its offsets modulo the original lengths. Thus the longest
cyclic overlap is the longest linear common substring of the doubled
words, capped at one shorter relator turn. Repeat with $D^{-1}$ to cover
the inverse sign.

For each unordered pair, it suffices to choose the shorter relation as
donor. If an overlap of length $k$ gives gain $2k-|D|$ with the longer
donor, it gives at least that gain with the shorter one. Hence this policy
is complete for the existence of an overlap satisfying $2k>|D|$ across
all distinct pairs and rotations. It returns the first pair and sign
with such an overlap, not the globally largest gain.

Absolute-generator counts give the additional bound
\[
 k\le M(R,D)=\sum_g\min(\#_gR,\#_gD).
\]
If $2M\le |D|$, the pair cannot help. Otherwise the LCS query may stop
at $M$: no one-turn cyclic overlap can exceed that value. This is a proved
upper bound, not a timeout heuristic. A direct prefix LCP sometimes
attains it before any overlap table is built. The generic LCS API likewise
accepts an explicit cap and returns the smaller of that cap and the exact
longest common substring length, together with a witness.

After rotating the relators to $R=PS$ and $D=PQ$, the producer checks the
prefix equality and replaces the target by the cyclic reduction of
$Q^{-1}S$. The donor is retained. This preserves normal closure, and
before further cancellation the target loses $2|P|-|D|>0$ letters.
Existing literal and compressed verifiers independently check the ordinary
\code{relator} move from the original knot presentation. The new matcher
is only a search mechanism, never a trusted shortcut in recognition.

The integration retains the fast paths: singleton elimination and
Whitehead shortening, the bounded explicit overlap bridge, then whole-donor
matching and quotient deletion. Only an oversized state that would otherwise
stall reaches the complete cyclic fallback. Every primitive shares the
arena's work allowance and cancellation callback. Node allocation, occurrence
table cells and cached overlap entries are bounded; resource exhaustion is
inconclusive and must not be mistaken for proof that no overlap exists.

\subsection{What complexity has actually been established}

Let $s$ be the input grammar size and $b$ the largest binary length of a
represented word length. There are $O(s^2)$ rule pairs and $O(b)$ dyadic
bands per required overlap set. Each band builds one polynomial-sized
occurrence table and makes a bounded number of exact compressed queries.
Each overlap progression needs at most six extensions. Every substring
used here is a slice of an input rule, its inverse, or a bounded binary
power of a slice. Slice height does not increase; powering adds $O(b)$
height and rules. The number and size of all intermediate grammars are
therefore polynomial in $s$ and $b$. Applying the established polynomial
bounds of the occurrence and equality primitives gives polynomial bit
complexity for one uncapped-resource LCS or cyclic-overlap query.

This is deliberately a local statement. It does not bound the number of
relator moves, the size of intermediate grammars across the entire search,
or whether a trivial knot group can always be simplified using strictly
shortening moves. The general subexponential and quasi-polynomial recognition
targets remain unproved. What has changed is that partial cyclic overlap
search is now implemented, rather than merely cited as a possible future
primitive.

\subsection{Validation and measurements}

The tests compare overlap sets for all 3844 ordered pairs of nonempty binary
words of length at most five, plus 300 signed random pairs with independently
randomized parses. Longest-common-substring tests cover 961 ordered binary
pairs including the empty word, 300 further signed random pairs and five
cap choices per pair, checking both lengths and witness positions. A
separate periodic test compares the critical-candidate maximum against every
member of the small progression, using independent noise on all four sides.
There are 350 randomized cyclic-pair trials checked against all literal
rotations and signs, including exact application and strict length decrease.

An exponential test represents $2^{500}$ overlap lengths by one progression
and finds a witness in $a^N b$ and $c a^N$ without expansion. Another test
starts with $a^N b^2,a^N c^2$: neither a Whitehead shortening nor a whole
donor deletion exists, but the new stage replaces the second word by
$b^{-2}c^2$. Total expanded relator length drops from $2N+4$ to $N+6$.
The supplied group remains inconclusive afterward; no knot verdict or
recognition-time speedup is inferred from this example. Invalid caps and
exhausted work and cache allowances are also checked.

The integrated suite passes 678 tests with Regina 7.4.1 in 163.634 seconds;
the targeted six-test run passes in 28.136 seconds.

\begin{center}\small
\begin{tabular}{@{}lrrr@{}}
\toprule
Scope & Before & A/A & Complete overlap \\
\midrule
Small 14, summed medians & 277.600 & 278.931 & 280.322 \\
Gordian, full query & 3425.051 & 3218.354 & 3264.758 \\
\bottomrule
\end{tabular}
\end{center}
Median milliseconds; all full knot queries complete.


All 225 full knot queries complete with identical certificate digests; none
reaches the new fallback. The small-fourteen sum of medians is 277.600,
278.931 and 280.322 milliseconds for archived, A/A and new search.
Gordian takes 3.425, 3.218 and 3.265 seconds respectively. The variation
between the historical controls prevents interpreting the lower Gordian
median as a demonstrated speedup. The audit supports preservation of the
existing successful paths.

\begin{center}\small
\begin{tabular}{@{}lrrrr@{}}
\toprule
LCS family & $h$ & Full query & Count bound & A/A \\
\midrule
Cyclic & 8 & 201.535 & 0.108 & 0.072 \\
Cyclic & 16 & 1969.533 & 0.164 & 0.105 \\
Cyclic & 32 & limit & 0.126 & 0.143 \\
Cyclic & 100 & limit & 0.462 & 0.320 \\
Cyclic & 500 & limit & 0.802 & 0.865 \\
Shifted & 8 & 0.671 & --- & 0.675 \\
Shifted & 16 & 1.252 & --- & 1.247 \\
Shifted & 32 & 2.359 & --- & 2.371 \\
Shifted & 100 & 7.285 & --- & 7.241 \\
Shifted & 500 & 37.239 & --- & 36.928 \\
\bottomrule
\end{tabular}
\end{center}
Median milliseconds including grammar construction. A/A repeats the bounded query for Cyclic and the full query for Shifted. Limits are censored queries.
\begin{center}\small
\begin{tabular}{@{}rrrr@{}}
\toprule
Partial move, $h$ & Milliseconds & Nodes & Charged work \\
\midrule
8 & 0.302 & 20 & 246 \\
32 & 0.615 & 44 & 606 \\
100 & 1.438 & 112 & 1626 \\
500 & 6.301 & 512 & 7626 \\
\bottomrule
\end{tabular}
\end{center}
Includes grammar construction, failed whole-donor search, complete cyclic search and checked application. This is not a knot verdict.


There are 125 measured LCS queries and twenty separate partial-move runs.
At $h=16$, the full cyclic-family query takes about 1.970 seconds and
1421272 work units; the count-bound query takes 0.164 milliseconds and
75 work units. At $h=32,100,500$, the full query exhausts its two-million-work
allowance. The bounded query still completes at $h=500$ in 0.802 milliseconds,
with 509 nonempty grammar nodes and 1527 work units. Submillisecond A/A
measurements are visibly noisy, so the stable work counts matter more than
precise ratios. Exhausted full queries are not completed-time denominators.

The shifted family tests a successful uncapped query that actually uses
periodic extension. At $h=500$, it takes 37.239 milliseconds versus
36.928 for A/A, with 3003 nodes and 27039 work units. The complete partial
move at the same exponent takes 6.301 milliseconds, 512 nodes and 7626
work units, including construction and checked application. These local
results do not demonstrate completion of the supplied group search or a
speedup on genuine knot recognition.


The audit uses archived search-loop code from \code{032ca00758} in the
common current recognition and replay harness, with an identical A/A arm.
There is one excluded warmup and five measured rounds per case, with
randomized arm order using seed 2810. Full queries include fresh planar-diagram
construction and independent replay. Their certificate digests are checked
outside the timer. All kernel timings include grammar construction; exhausted
queries remain censored. The cyclic query family compares a full LCS query
with the proved shared-count cap on $(a^N b^2)^2,(a^N c^2)^2$, whose exact
LCS is $N$. A second family uses $a^N b,c a^N$ and compares two identical
uncapped queries. A separate move audit includes failed whole-donor search,
complete cyclic search and checked application.

Whole recognition allows 18 seconds, with 15 seconds and twenty million
work units for its group stage. Kernels allow two million work units and
100000 nodes/table cells. Raw results are in
\path{results/compressed_lcs_20261008.json}.
Test and build logs are in \path{synthesis/data/compressed-lcs-*}.

No new incoming report was present; the suggested \code{openai/math}
head remained at the previously recorded \code{fd4aeeb2ee}. This iteration implements
and checks the previously identified compressed-LCS research direction.
